\chapter{الحذف \texorpdfstring{$\beta$}{بيتا}}\label{ch:b1:elimination}

عالج 2-برومو-2-ميثيل البروبان بإيثوكسيد الصوديوم في الإيثانول عند درجة حرارة
الغرفة تجد أن الناتج الرئيسي إيثر؛ وسخّن الخليط نفسه تجد أن
الناتج الرئيسي ألكين، هو 2-ميثيل البروبين، لا مجموعة إيثوكسي
فيه. فأيون الإيثوكسيد لم يعمل \omterm{def:b1:electronic-effects:nucleophile}{محبًا للنواة} يهاجم الكربون، بل
قاعدةً، تنزع بروتونًا من الكربون المجاور بينما يغادر
البروميد. وهذا التنافس بين الاستبدال والحذف يسري في
التخليق العضوي كله. يصف هذا الفصل آليتي الحذف،
وشرطهما الهندسي اللافت، والقاعدة التي
تتنبأ بالألكين المتكوّن، وكيف يُوجَّه تفاعل في هذا الاتجاه أو
ذاك.

\begin{recall}
SN1 وSN2 (\cref{ch:b1:nucleophilic-substitution})؛ وإسقاطات نيومان،
وتشكّلات الهكسان الحلقي والواصفان $Z$/$E$
(\cref{ch:b1:stereochemistry-in-depth})؛ \omterm{def:b1:electronic-effects:intermediates}{والكاتيونات الكربونية} والقواعد
(\cref{ch:b1:electronic-effects}).
\end{recall}

\section{الآلية E2}

\begin{definition}[الحذف \texorpdfstring{$\beta$}{بيتا}، E2 وE1]\label{def:b1:elimination:elimination}
\emph{الحذف $\beta$}\index{حذف $\beta$} ينزع
\omterm{def:b1:electronic-effects:nucleophile}{مجموعة مغادرة} Y من كربون (الكربون $\alpha$) وبروتونًا من
كربون مجاور (كربون $\beta$)، فتتكوّن رابطة ثنائية بينهما.
وفي \emph{الآلية E2}\index{آلية E2} تأخذ القاعدة
البروتون بينما تغادر Y، في \omterm{def:b1:elementary-steps:elementary-step}{خطوة أولية} واحدة؛ وفي \emph{الآلية
E1}\index{آلية E1} تغادر Y أولًا، فيتكوّن \omterm{def:b1:electronic-effects:intermediates}{كاتيون كربوني}،
يفقد بعد ذلك بروتونًا $\beta$.
\end{definition}

\begin{proposition}[قانون سرعة E2]\label{prop:b1:elimination:e2-law}
من أجل E2، $v = k[\mathrm{RY}][\mathrm{B}]$، حيث B القاعدة.
\end{proposition}

\begin{proof}
\omterm{def:b1:elementary-steps:elementary-step}{خطوة أولية} واحدة ثنائية الجزيئية، تشارك فيها \omterm{def:b1:nucleophilic-substitution:substitution}{الركيزة} والقاعدة: قانون
فعل الكتلة \omterm{def:b1:elementary-steps:elementary-step}{للخطوات الأولية} (\cref{ch:b1:elementary-steps}).
\end{proof}

\begin{definition}[الوضع المضاد المستوي والوضع المتزامن المستوي]\label{def:b1:elimination:periplanar}
تكون رابطتان على ذرتين متجاورتين \emph{في وضع مضاد مستوٍ}\index{وضع مضاد مستو}
حين تقعان في مستوٍ واحد على جهتين متقابلتين (\omterm{def:b1:stereochemistry-in-depth:conformations}{زاوية الالتواء} $180^\circ$)،
و\emph{في وضع متزامن مستوٍ}\index{وضع متزامن مستو} حين تقعان في مستوٍ واحد في
الجهة نفسها ($0^\circ$).
\end{definition}

\begin{omfigure}
\begin{tikzpicture}[font=\small, >={Stealth}]
  \pic (n) at (0,0) {omnewman={90}{270}};
  \node at (n-f1) {H}; \node at (n-f2) {\ce{CH3}}; \node at (n-f3) {H};
  \node at (n-b1) {Br}; \node at (n-b2) {\ce{CH3}}; \node at (n-b3) {H};
  \node (b) at (-2.6,2.0) {\ce{EtO-}};
  \draw[->, thick, omThm] (b) to[bend left=20] ($(n-f1)+(-0.15,0.1)$);
  \draw[->, thick, omThm] (0.12,0.8) to[bend left=40] (0.12,0.25);
  \draw[->, thick, omThm] (0.12,-0.6) to[bend right=35] ($(n-b1)+(0.25,0.15)$);
  \node[below] at (0,-2.0) {\foreignlanguage{arabic}{\babelsublr{H} و\babelsublr{Br} في وضع مضاد مستوٍ}};
  \draw[->, thick] (2.4,0) -- (3.6,0) node[midway, above] {E2};
  \node at (5.6,0) {\chemfig{C(-[:150]H_3C)(-[:210]H)=C(-[:30]H)(-[:-30]CH_3)}};
  \node at (5.6,-0.9) {\foreignlanguage{arabic}{بوت-\babelsublr{2}-ين-\babelsublr{\cip{E}}}};
\end{tikzpicture}

{\small التفاعل E2 للمركب 2-بروموبوتان، منظورًا إليه على امتداد الرابطة C3--C2 (C3 في
الأمام، وC2 الحامل لذرة Br في الخلف) في واحد من تشكّليه المضادين
المستويين؛ \omterm{def:b1:stereochemistry-in-depth:isomers}{والتشكّل} الآخر، الذي يكون فيه الهيدروجين الآخر للكربون C3 مضادًا لذرة Br،
يعطي بوت-2-ين-\cip{Z}. وتتحرك ثلاثة أزواج
إلكترونية معًا: تأخذ القاعدة البروتون، ويصير زوج C--H
الرابطة $\pi$، ويغادر زوج C--Br مع البروم. والمجموعات المتقابلة
في \omterm{def:b1:stereochemistry-in-depth:representations}{إسقاط نيومان} تنتهي في الجهة نفسها من الرابطة الثنائية.}
\end{omfigure}

\begin{proposition}[الآلية E2 مضادة ونوعية فراغيًا]\label{prop:b1:elimination:anti}
تجري E2 انطلاقًا من \omterm{def:b1:stereochemistry-in-depth:isomers}{التشكّل} الذي تكون فيه الرابطتان C--H وC--Y
\omterm{def:b1:elimination:periplanar}{في وضع مضاد مستوٍ}. وهي نوعية فراغيًا: فالركائز المتماكبة لامرآويًا تعطي
\omterm{def:b1:stereochemistry-in-depth:isomers}{متماكبات فراغية} مختلفة من الألكين، يحددها الترتيب المضاد.
\end{proposition}

\begin{proof}
في \omterm{def:b1:elementary-steps:energy-profile}{الحالة الانتقالية} يجب أن يدخل \omterm{def:b1:lewis-resonance-vsepr:lewis}{الزوج الرابط} C--H المنطقة
المضادة للترابط للرابطة C--Y بينما يصير الكربونان
مثلثيين؛ ويكون التداخل أفضل ما يكون حين تكون الرابطتان متوازيتين، في
مستوٍ واحد. والترتيب المضاد متخالف (طاقته منخفضة) ويُبقي القاعدة
بعيدة عن \omterm{def:b1:electronic-effects:nucleophile}{المجموعة المغادرة}؛ أما المتزامن فمتطابق. ومع H وY متضادتين،
تتحدد المجموعتان الأخريان على كل كربون: فاللتان تقعان في الجهة
نفسها في \omterm{def:b1:stereochemistry-in-depth:representations}{إسقاط نيومان} تنتهيان في وضع \textit{cis} على الرابطة الثنائية.
\omterm{def:b1:stereochemistry-in-depth:enantiomers}{والمتماكب اللامرآوي} فيه مجموعتان متبادلتان على كربون واحد، فيعطي
الألكين الآخر.
\end{proof}

\begin{proposition}[الآلية E2 على هكسان حلقي]\label{prop:b1:elimination:cyclohexane-e2}
على حلقة الهكسان الحلقي، تتطلب E2 أن تكون \omterm{def:b1:electronic-effects:nucleophile}{المجموعة المغادرة} وهيدروجين $\beta$
محوريين كلاهما (\textit{trans}-ثنائي المحور)؛ \omterm{def:b1:electronic-effects:nucleophile}{والمجموعة المغادرة} المثبتة
في موضع \omterm{def:b1:stereochemistry-in-depth:conformations}{استوائي} لا تستطيع التفاعل حتى تنقلب الحلقة.
\end{proposition}

\begin{proof}
على كربونين متجاورين في \omterm{def:b1:stereochemistry-in-depth:conformations}{كرسي}، تكون رابطتان محوريتان \omterm{def:b1:elimination:periplanar}{في وضع مضاد مستوٍ} (إحداهما
إلى الأعلى والأخرى إلى الأسفل، موازيتين للمحور)؛ أما الرابطة \omterm{def:b1:stereochemistry-in-depth:conformations}{الاستوائية} فمضادة لروابط
C--C في الحلقة فقط، ولا تكون أبدًا مضادة لرابطة C--H.
\end{proof}

\begin{omfigure}
\begin{tikzpicture}[font=\small]
  \pic (a) at (0,0) {omchair};
  \draw[thick, omThm] (a-c1) -- (a-a1) node[above] {Cl};
  \draw[thick] (a-c1) -- (a-e1) node[right] {H};
  \draw[thick, omDef] (a-c2) -- (a-a2) node[below] {H};
  \draw[thick, omDef] (a-c6) -- (a-a6) node[below] {H};
  \node[align=left, anchor=west] at (3.0,0.2) {\foreignlanguage{arabic}{\babelsublr{Cl} محورية إلى الأعلى على \babelsublr{C1}؛}\\\foreignlanguage{arabic}{\babelsublr{H} محوري إلى الأسفل على الجارين كليهما:}\\\foreignlanguage{arabic}{ذرتا \babelsublr{H} في وضع مضاد مستوٍ}};
\end{tikzpicture}

{\small الترتيب \textit{trans}-ثنائي المحور على \omterm{def:b1:stereochemistry-in-depth:conformations}{كرسي}: الكلور \omterm{def:b1:stereochemistry-in-depth:conformations}{المحوري}
\omterm{def:b1:elimination:periplanar}{في وضع مضاد مستوٍ} بالنسبة إلى الهيدروجينين المحوريين على الكربونين المجاورين
(بالأزرق).}
\end{omfigure}

\section{الآلية E1}

\begin{proposition}[قانون سرعة E1]\label{prop:b1:elimination:e1-law}
من أجل E1، $v = k[\mathrm{RY}]$، والسرعة هي سرعة التأين.
\end{proposition}

\begin{proof}
كما في SN1 (\cref{prop:b1:nucleophilic-substitution:sn1-law}): يتكوّن
\omterm{def:b1:electronic-effects:intermediates}{الكاتيون الكربوني} في خطوة بطيئة ويُستهلك بسرعة، هنا بفقد
بروتون لصالح المذيب أو \omterm{def:b1:predominant-reaction:strong-acid}{قاعدة ضعيفة}؛ ويُبقي \omterm{def:b1:elementary-steps:qssa}{تقريب الحالة المستقرة}
$v = k_1[\mathrm{RY}]$.
\end{proof}

تشترك E1 في كاتيونها الكربوني مع SN1: فالركائز الثالثية في \omterm{def:b1:intermolecular-forces-solvents:solvent-classes}{المذيبات
البروتونية} تعطي الاثنين، والتسخين يحابي الألكين. ولأن
\omterm{def:b1:electronic-effects:intermediates}{الكاتيون الكربوني} مستوٍ والبروتون يُفقد لاحقًا، فليس \omterm{def:b1:elimination:elimination}{للآلية E1}
شرط هندسي وهي ليست نوعية فراغيًا؛ وتعطي أساسًا الألكين الأكثر
ثباتًا، وهو عادةً الألكين \cip{E}.

\begin{omfigure}
\setchemfig{atom sep=2.1em}
\schemestart
\chemfig{C(-[2]CH_3)(-[4]H_3C)(-[6]CH_3)-Br}
\arrow{->[بطيئة]}
\chemfig{C^{+}(-[2]CH_3)(-[:210]H_3C)(-[:-30]CH_3)}
\arrow{->[سريعة][\ce{-H+}]}
\chemfig{C(-[2]CH_3)(-[:210]H_3C)=[:-30]CH_2}
\schemestop
\par\vspace{4pt}

{\small E1: تأين 2-برومو-2-ميثيل البروبان إلى \omterm{def:b1:electronic-effects:intermediates}{الكاتيون الكربوني}، ثم
فقد بروتون من مجموعة ميثيل لصالح المذيب، فيتكوّن
2-ميثيل البروبين.}
\end{omfigure}

\section{الانتقائية الموضعية: قاعدة زايتسيف}

\begin{definition}[التفاعلات الانتقائية موضعيًا والانتقائية فراغيًا]\label{def:b1:elimination:selectivity}
التفاعل الذي يمكن أن يعطي عدة \omterm{def:b1:stereochemistry-in-depth:isomers}{متماكبات بنيوية} لكنه يكوّن واحدًا منها
أساسًا \emph{تفاعل انتقائي موضعيًا}\index{تفاعل انتقائي موضعيًا}؛
والذي يمكن أن يعطي عدة \omterm{def:b1:stereochemistry-in-depth:isomers}{متماكبات فراغية} لكنه يكوّن واحدًا منها أساسًا
\emph{تفاعل انتقائي فراغيًا}\index{تفاعل انتقائي فراغيًا}.
\end{definition}

\begin{proposition}[قاعدة زايتسيف]\label{prop:b1:elimination:zaitsev}
حين يمكن نزع عدة هيدروجينات $\beta$، يكون الألكين الرئيسي
عادةً الأكثر استبدالًا (الأكثر ثباتًا)، والمتماكب \cip{E}
لا \cip{Z} (\emph{قاعدة زايتسيف}\index{قاعدة زايتسيف}).
ويمكن للقواعد الضخمة، وللقيود الهندسية مثل شرط الوضع ثنائي المحور
\textit{trans}، أن تتغلب عليها.
\end{proposition}

\begin{proof}
مجموعات الألكيل على كربوني رابطة ثنائية تثبّتها (كما
تثبّت \omterm{def:b1:electronic-effects:intermediates}{الكاتيونات الكربونية})، والألكينات \cip{E} تتجنب ازدحام
المجموعات \textit{cis}. وللحالات الانتقالية في E2 وE1 طابع جزئي من
الرابطة الثنائية، فيتكوّن الألكين الأكثر ثباتًا أسرع.
والقاعدة الضخمة تبلغ الهيدروجينات الأقل إعاقةً بسهولة أكبر؛
والهيدروجين الذي لا يمكن أن يكون \omterm{def:b1:elimination:periplanar}{في وضع مضاد مستوٍ} بالنسبة إلى \omterm{def:b1:electronic-effects:nucleophile}{المجموعة المغادرة} لا
تنزعه E2 أبدًا، مهما كان ثبات الألكين الذي كان سيعطيه.
\end{proof}

\begin{method}[التنبؤ بالألكين]\label{met:b1:elimination:predict-alkene}
\begin{enumerate}
  \item اذكر كربونات $\beta$ التي تحمل هيدروجينات.
  \item من أجل E2، لا تُبقِ إلا الهيدروجينات التي يمكن أن تكون \omterm{def:b1:elimination:periplanar}{في وضع مضاد مستوٍ} بالنسبة إلى Y (على
        حلقة: ثنائية المحور \textit{trans}، في \omterm{def:b1:stereochemistry-in-depth:conformations}{كرسي} يمكن أن يتكوّن فعلًا).
  \item ومن بينها، فضّل الهيدروجين الذي يعطي الألكين الأكثر استبدالًا
        (زايتسيف)، ما لم تكن القاعدة ضخمة.
  \item ولكل منها، ارسم \omterm{def:b1:stereochemistry-in-depth:representations}{إسقاط نيومان} في \omterm{def:b1:stereochemistry-in-depth:conformations}{التشكّل المضاد} لتقرأ
        الهندسة \cip{E}/\cip{Z}.
\end{enumerate}
\end{method}

\section{استبدال أم حذف؟}

\begin{proposition}[الاستبدال مقابل الحذف]\label{prop:b1:elimination:sn-vs-e}
يحابي الحذفَ القواعدُ القوية الضخمة، ودرجةُ الحرارة المرتفعة
والركائزُ الثالثية أو المزدحمة؛ ويحابي الاستبدالَ المحباتُ للنواة الجيدة التي هي
\omterm{def:b1:predominant-reaction:strong-acid}{قواعد ضعيفة}، ودرجةُ الحرارة المنخفضة والركائزُ الأولية.
\end{proposition}

\begin{proof}
قوة القاعدة تحابي الهجوم على الهيدروجين؛ والضخامة تعيق الهجوم على
كربون مزدحم لا على هيدروجين مكشوف. \omterm{def:b1:electronic-effects:carbon-class}{والكربون الثالثي}
لا تبلغه SN2، وكاتيونه الكربوني يفقد البروتونات بسهولة. والحذف
يحوّل جزيئًا واحدًا إلى جزيئين أو ثلاثة (ألكين، \omterm{def:b1:electronic-effects:nucleophile}{ومجموعة مغادرة}، وقاعدة
متبرتنة): فهو محابى كلما ارتفعت درجة الحرارة، وهي نتيجة يسوّغها
مجلد السنة الثانية.
\end{proof}

\begin{omfigure}
\begin{tabular}{llll}
\hline
\omterm{def:b1:nucleophilic-substitution:substitution}{الركيزة} & \omterm{def:b1:predominant-reaction:strong-acid}{قاعدة قوية}، صغيرة & \omterm{def:b1:predominant-reaction:strong-acid}{قاعدة قوية}، ضخمة & \omterm{def:b1:predominant-reaction:strong-acid}{قاعدة ضعيفة}، \omterm{def:b1:electronic-effects:nucleophile}{محب للنواة} جيد \\
\hline
أولية & SN2 (وقليل من E2) & E2 & SN2 \\
ثانوية & E2 (وSN2) & E2 & SN2 (لابروتوني)، SN1/E1 (بروتوني) \\
ثالثية & E2 & E2 & SN1/E1 (بروتوني) \\
\hline
\end{tabular}

\smallskip
{\small توجيه أول: التفاعل الرئيسي في كل حالة. والتسخين
يحابي الحذف عمودًا بعد عمود.}
\end{omfigure}

\section{تمارين}

\begin{exercise}[$\star$]\label{exo:b1:elimination:1}
أعطِ الألكينات التي يمكن أن تتكوّن بحذف HBr من 2-بروموبوتان،
ومن 2-برومو-2-ميثيل البوتان ومن البروموهكسان الحلقي، والألكين الذي تتنبأ \omterm{prop:b1:elimination:zaitsev}{قاعدة زايتسيف}
بأنه الرئيسي.
\end{exercise}

\begin{exercise}[$\star$]\label{exo:b1:elimination:2}
اكتب \omterm{def:b1:rate-laws:order}{قانوني سرعة} E1 وE2. وكيف تميّز بينهما
تجريبيًا في حالة 2-برومو-2-ميثيل البروبان في إيثانول يحتوي إيثوكسيدًا؟
\end{exercise}

\begin{exercise}[$\star$]\label{exo:b1:elimination:3}
ارسم إسقاطات نيومان للمركب 2-بروموبوتان على امتداد C2--C3 في تشكّلاته المتخالفة
الثلاثة، وبيّن أيها يمكن أن يخضع \omterm{def:b1:elimination:elimination}{للآلية E2}.
\end{exercise}

\begin{exercise}[$\star$]\label{exo:b1:elimination:4}
توقّع التفاعل الرئيسي (SN1، SN2، E1، E2): البروموإيثان مع هيدروكسيد
الصوديوم في الماء عند \qty{25}{\celsius}؛ و2-برومو-2-ميثيل البروبان مع
ثالثي بوتوكسيد البوتاسيوم؛ و2-بروموبروبان مع يوديد الصوديوم في البروبانون؛
و2-برومو-2-ميثيل البروبان في إيثانول ساخن.
\end{exercise}

\begin{exercise}[$\star\star$]\label{exo:b1:elimination:5}
بيّن، بإسقاطات نيومان، أن الحذف E2 للمركب HBr من
\omterm{def:b1:stereochemistry-in-depth:enantiomers}{المتماكبين اللامرآويين} للمركب 2-برومو-3-ميثيل البنتان اللذين يمكن أن يفقدا هيدروجين C3
يعطي \omterm{def:b1:stereochemistry-in-depth:isomers}{المتماكبين الفراغيين} للمركب 3-ميثيل بنت-2-ين.
\end{exercise}

\begin{exercise}[$\star\star$]\label{exo:b1:elimination:6}
اكتب آلية نزع الماء من 2-ميثيل بوتان-2-أول في حمض الكبريتيك
المركّز الساخن (E1) وأعطِ الألكين الرئيسي.
\end{exercise}

\begin{exercise}[$\star\star$]\label{exo:b1:elimination:7}
فسّر لماذا يتفاعل 1-برومو-4-ثالثي بوتيل الهكسان الحلقي-\textit{trans} ببطء
شديد \omterm{def:b1:elimination:elimination}{بالآلية E2} بينما يتفاعل متماكبه \textit{cis} بسرعة. (تبقى مجموعة ثالثي البوتيل
\omterm{def:b1:stereochemistry-in-depth:conformations}{استوائية}.)
\end{exercise}

\begin{exercise}[$\star\star$]\label{exo:b1:elimination:8}
يعطي 2-برومو-2-ميثيل البوتان مع إيثوكسيد الصوديوم أساسًا
2-ميثيل بوت-2-ين؛ ومع ثالثي بوتوكسيد البوتاسيوم، أساسًا
2-ميثيل بوت-1-ين. فسّر.
\end{exercise}

\begin{exercise}[$\star\star$]\label{exo:b1:elimination:9}
لماذا لا يُرصد الحذف أبدًا مع البروموميثان؟ ولا مع
1-برومو-2,2-ثنائي ميثيل البروبان \omterm{def:b1:elimination:elimination}{بالآلية E2}؟
\end{exercise}

\begin{exercise}[$\star\star\star$]\label{exo:b1:elimination:10}
في الإيثانول عند \qty{25}{\celsius}، يعطي 2-برومو-2-ميثيل البروبان خليطًا
من الإيثر (SN1) والألكين (E1) لا تتعلق نسبتاهما
بتركيز \omterm{def:b1:nucleophilic-substitution:substitution}{الركيزة}. بيّن أن الناتجين يأتيان من
\omterm{def:b1:electronic-effects:intermediates}{الكاتيون الكربوني} نفسه وأن نسبتهما هي نسبة ثابتي
سرعة.
\end{exercise}

\begin{exercise}[$\star\star\star$]\label{exo:b1:elimination:11}
يعطي أحد \omterm{def:b1:stereochemistry-in-depth:enantiomers}{المتماكبين اللامرآويين} للمركب 1,2-ثنائي برومو-1,2-ثنائي فينيل الإيثان (شكل الميزو)،
\omterm{def:b1:elimination:elimination}{بالآلية E2} مع الإيثوكسيد، \omterm{def:b1:stereochemistry-in-depth:isomers}{متماكبًا فراغيًا} وحيدًا من 1-برومو-1,2-ثنائي فينيل الإيثين.
أوجده، \omterm{def:b1:stereochemistry-in-depth:representations}{بإسقاط نيومان}.
\end{exercise}

\begin{exercise}[$\star\star\star$]\label{exo:b1:elimination:12}
مستعينًا بطاقات \cref{ch:b1:stereochemistry-in-depth}، فسّر لماذا
يمكن أن يتباطأ تفاعل E2 على هكسان حلقي بعامل يبلغ المئات
حين يجب أن تصير \omterm{def:b1:electronic-effects:nucleophile}{المجموعة المغادرة} \omterm{def:b1:stereochemistry-in-depth:conformations}{محورية} في \omterm{def:b1:stereochemistry-in-depth:conformations}{كرسي} مزدحم. وعبّر عن
السرعة بجداء كسر من المتشكّلات \omterm{def:b1:rate-laws:order}{وثابت سرعة}.
\end{exercise}

\section{مسألة: كلوريدا المنثيل والنيومنثيل}

\begin{problem}[{مسألة نهاية الأسبوع --- كرسيا كلوريدين متماكبين لامرآويًا،
والهيدروجينات المتاحة في وضع مضاد مستوٍ في كل منهما، والألكينات
المتكوّنة، ونسبة سرعتي E2 لهما}]
\label{pb:b1:elimination:1}
كلوريد المنثيل وكلوريد النيومنثيل هما كلوريدا المنثول
والنيومنثول (\cref{ch:b1:stereochemistry-in-depth}): على الحلقة، يحمل C1
ذرة Cl، وC2 مجموعة الإيزوبروبيل، وC5 مجموعة الميثيل. وفي كلوريد
المنثيل يمكن أن تكون المجموعات الثلاث كلها \omterm{def:b1:stereochemistry-in-depth:conformations}{استوائية}؛ وفي كلوريد النيومنثيل،
تكون Cl \omterm{def:b1:stereochemistry-in-depth:conformations}{محورية} حين تكون مجموعتا الألكيل استوائيتين. ويُعالج الاثنان
بإيثوكسيد الصوديوم في الإيثانول، وهذا تفاعل E2. معطيات المسألة:
\qty{95}{\percent} من كلوريد النيومنثيل في \omterm{def:b1:stereochemistry-in-depth:conformations}{الكرسي} الذي تكون فيه Cl \omterm{def:b1:stereochemistry-in-depth:conformations}{محورية}؛
ومن أجل كلوريد المنثيل، يمثّل \omterm{def:b1:stereochemistry-in-depth:conformations}{الكرسي} الذي تكون فيه Cl \omterm{def:b1:stereochemistry-in-depth:conformations}{محورية} (والمجموعات الثلاث كلها \omterm{def:b1:stereochemistry-in-depth:conformations}{محورية})
كسرًا قدره \num{5.0e-3}؛ وافترض أن كل هيدروجين \omterm{def:b1:elimination:periplanar}{في وضع مضاد مستوٍ} في
\omterm{def:b1:stereochemistry-in-depth:conformations}{كرسي} نشط يتفاعل \omterm{def:b1:rate-laws:order}{بثابت السرعة} نفسه. ويعطي كلوريد النيومنثيل
\qty{78}{\percent} من الألكين الثلاثي الاستبدال.

\medskip
\noindent\textbf{الجزء الأول --- الكرسيان.}
\begin{enumerate}
  \item ارسم \omterm{def:b1:stereochemistry-in-depth:conformations}{الكرسي} المفضل لكلوريد المنثيل.
  \item ارسم \omterm{def:b1:stereochemistry-in-depth:conformations}{الكرسي} المفضل لكلوريد النيومنثيل.
  \item هل الكلوريدان \omterm{def:b1:stereochemistry-in-depth:enantiomers}{متماكبان مرآويان} أم \omterm{def:b1:stereochemistry-in-depth:enantiomers}{متماكبان لامرآويان}؟
  \item اذكر الشرط الهندسي \omterm{def:b1:elimination:elimination}{للآلية E2} على حلقة.
  \item في أي \omterm{def:b1:stereochemistry-in-depth:conformations}{كرسي} يمكن لكلوريد المنثيل أن يتفاعل؟ ارسمه.
  \item لماذا كان ذلك \omterm{def:b1:stereochemistry-in-depth:conformations}{الكرسي} نادرًا؟ استعمل كلفة الطاقة للمجموعات \omterm{def:b1:stereochemistry-in-depth:conformations}{المحورية}.
\end{enumerate}

\medskip
\noindent\textbf{الجزء الثاني --- الهيدروجينات المضادة.}
\begin{enumerate}[resume]
  \item اذكر كربونات $\beta$ في الكلوريد (C2 وC6)
        وهيدروجيناتها.
  \item في \omterm{def:b1:stereochemistry-in-depth:conformations}{الكرسي} النشط لكلوريد النيومنثيل، أي الهيدروجينات تكون
        \omterm{def:b1:elimination:periplanar}{في وضع مضاد مستوٍ} بالنسبة إلى Cl؟
  \item في \omterm{def:b1:stereochemistry-in-depth:conformations}{الكرسي} النشط لكلوريد المنثيل، هل هيدروجين C2
        \omterm{def:b1:stereochemistry-in-depth:conformations}{محوري}؟ وهل أحد هيدروجيني C6 \omterm{def:b1:stereochemistry-in-depth:conformations}{محوري}؟
  \item كم هيدروجينًا \omterm{def:b1:elimination:periplanar}{في وضع مضاد مستوٍ} يقدّم كل كلوريد؟
  \item ارسم \omterm{def:b1:stereochemistry-in-depth:representations}{إسقاط نيومان} على امتداد C1--C6 \omterm{def:b1:stereochemistry-in-depth:conformations}{للكرسي} النشط
        لكلوريد المنثيل.
  \item لماذا لا يستطيع حذف \omterm{def:b1:elimination:periplanar}{في وضع متزامن مستوٍ} أن ينقذ
        التفاعل؟
\end{enumerate}

\medskip
\noindent\textbf{الجزء الثالث --- النواتج.}
\begin{enumerate}[resume]
  \item سمِّ الألكين المتكوّن بنزع هيدروجين C2 والألكين
        المتكوّن بنزع هيدروجين من C6.
  \item أيهما تحابي \omterm{prop:b1:elimination:zaitsev}{قاعدة زايتسيف}، ولماذا؟
  \item ماذا يعطي كلوريد النيومنثيل؟ وهل النسبة 78/22 متسقة
        مع \omterm{prop:b1:elimination:zaitsev}{قاعدة زايتسيف}؟
  \item ماذا يعطي كلوريد المنثيل؟ وهل هذا متسق مع
        \omterm{prop:b1:elimination:zaitsev}{قاعدة زايتسيف}؟
  \item أي التفاعلين نوعي فراغيًا، أو انتقائي موضعيًا، أو كلاهما؟
  \item لماذا كان الاستبدال (SN2) ثانويًا مع الإيثوكسيد هنا؟
\end{enumerate}

\medskip
\noindent\textbf{الجزء الرابع --- السرعتان.}
\begin{enumerate}[resume]
  \item اكتب سرعة كل تفاعل دالةً في كسر
        \omterm{def:b1:stereochemistry-in-depth:conformations}{الكرسي} النشط وعدد الهيدروجينات \omterm{def:b1:elimination:periplanar}{في وضع مضاد مستوٍ}.
  \item احسب السرعة النسبية لكلوريد النيومنثيل.
  \item احسب السرعة النسبية لكلوريد المنثيل.
  \item احسب نسبة السرعتين $k(\text{نيومنثيل})/k(\text{منثيل})$.
\end{enumerate}
\end{problem}
